\chapter{Teori Belajar Statistik}
\label{ch:teori}

Kuliah Statistical Learning Theory and Applications dari Werner Zellinger saya ikuti pada musim panas 2025. Naskahnya pendek,
bukunya tebal: rujukan utamanya adalah \emph{Learning Theory from First Principles} \citep{bach2024}. Pertanyaan yang diajukan
di halaman pertama naskah sederhana sekali. Kalau data tak terhingga, model apa yang terbaik, dan berapa galat terbaiknya?
Kapan sebuah algoritme belajar dijamin semakin baik dengan bertambahnya data? Dan apakah belajar mungkin tanpa asumsi apa pun
tentang data? Bab ini mengikuti jawaban-jawabannya. Ia lebih formal dari \cref{ch:data}, tetapi setiap definisi punya alasan
yang konkret.

\section{Merumuskan belajar}

Data adalah $n$ pasangan $(\vx_i,y_i)$ yang ditarik independen dari distribusi $P$ yang tidak diketahui. Algoritme belajar adalah
fungsi yang memetakan himpunan data ke sebuah model $f$. Loss $\ell(z,y)$ mengukur kerugian memakai prediksi $z$ ketika jawabannya
$y$. Risiko model adalah loss harapan atas distribusi data, $R(f)=\E_{(\vx,y)\sim P}[\ell(f(\vx),y)]$, dan risiko empiris adalah rata-ratanya
di data, $\hat R(f)=\tfrac1n\sum_i\ell(f(\vx_i),y_i)$.

Model terbaik yang mungkin bisa ditulis tanpa melihat data sama sekali. Karena $R(f)=\E_{\vx}\big[\E[\ell(f(\vx),y)\mid\vx]\big]$, risiko
diminimalkan dengan memilih, untuk setiap $\vx$ secara terpisah, prediksi $z$ yang meminimalkan loss harapan bersyarat. Fungsi
yang melakukannya disebut \istilah{risk minimizer} $f^*$, dan risikonya, $R^*$, adalah risiko minimal. Untuk regresi dengan loss
kuadrat, $\E[(z-y)^2\mid\vx]$ adalah parabola dalam $z$ yang minimum di rata-rata bersyarat, sehingga $f^*(\vx)=\E[y\mid\vx]$, sama dengan
yang kita turunkan untuk satu konstanta di \cref{ch:stat}. Untuk klasifikasi biner dengan loss 0--1, prediksi $z=1$ salah dengan
peluang $P(y=0\mid\vx)$ dan prediksi $z=0$ salah dengan peluang $P(y=1\mid\vx)$, sehingga yang terbaik adalah memilih kelas dengan
peluang bersyarat lebih besar, dan
\begin{equation}
R^*=\E_{\vx}\Big[\min\big\{P(y=1\mid\vx),\ 1-P(y=1\mid\vx)\big\}\Big].
\end{equation}
Itulah galat Bayes dari \cref{ch:stat}, dalam bahasa yang lebih umum. Kalau $P$ diketahui, belajar menjadi trivial. Seluruh
kesulitan berasal dari kenyataan bahwa $P$ hanya diketahui lewat sampel.

\section{Dua sumber galat}

Algoritme memilih model dari sebuah kelas $\mathcal F$, misalnya semua fungsi linear atau semua jaringan berukuran tertentu.
Galat berlebih model yang dipilih, $\hat f$, terhadap yang terbaik yang mungkin bisa dipecah menjadi dua:
\begin{equation}
R(\hat f)-R^*=\underbrace{\Big(R(\hat f)-\inf_{f\in\mathcal F}R(f)\Big)}_{\text{galat estimasi}}
+\underbrace{\Big(\inf_{f\in\mathcal F}R(f)-R^*\Big)}_{\text{galat aproksimasi}}.
\end{equation}
Galat aproksimasi adalah harga memilih kelas yang terlalu sempit: kalau fungsi sebenarnya melengkung dan kita hanya mengizinkan
garis lurus, tidak ada jumlah data yang bisa menolong. Galat ini tidak bergantung pada data. Galat estimasi adalah harga memilih
dari data yang terbatas: semakin kaya kelasnya, semakin mudah kita memilih anggota yang kebetulan cocok dengan noise di data latih.
Galat ini mengecil dengan bertambahnya data dan membesar dengan bertambahnya kapasitas kelas. Kelas yang lebih kaya menurunkan
yang pertama dan menaikkan yang kedua, bentuk umum dari tawar-menawar bias dan varians di \cref{ch:data}. Kapasitas dikendalikan
dengan mengurangi parameter atau dengan menambahkan penalti norma, seperti weight decay.

\section{Kualitas sebuah algoritme}

Kuliah membedakan beberapa cara menilai algoritme. Algoritme \istilah{konsisten} kalau risiko harapan model yang dihasilkannya
mendekati $R^*$ ketika $n$ menuju tak hingga. Versi PAC, \emph{probably approximately correct}, meminta lebih: dengan peluang
paling sedikit $1-\delta_n$, galat berlebihnya tidak lebih dari $\epsilon_n$, dengan kedua batas mengecil ke nol. Algoritme
\istilah{universally consistent} konsisten untuk setiap distribusi dalam sebuah kelas yang besar. Risiko \istilah{minimax},
\begin{equation}
\inf_{\mathcal A}\ \sup_{P\in\mathcal P}\ \E_{\text{data}\sim P^n}\big[R(\mathcal A(\text{data}))\big]-R^*,
\end{equation}
menilai algoritme terbaik pada distribusi terburuknya. Salah satu tujuan utama teori belajar adalah algoritme yang mencapai laju
penurunan risiko minimax yang optimal.

\section{Tidak ada makan siang gratis}

Teorema \emph{no free lunch} dalam kuliah menjawab pertanyaan ketiga di awal bab. Untuk klasifikasi biner dengan loss 0--1, kalau
kelas distribusi yang diizinkan adalah \emph{semua} distribusi pada ruang input yang tak terhingga, maka untuk setiap algoritme dan
setiap ukuran data, ada distribusi yang membuat galat berlebih harapannya paling sedikit $\tfrac12$. Gagasan pembuktiannya bisa
diceritakan tanpa rumus. Pilih distribusi yang tersebar di banyak sekali titik, sehingga $n$ contoh hampir pasti tidak pernah
menyentuh titik-titik uji. Lalu beri label pada titik-titik yang tidak terlihat itu secara acak. Apa pun yang dilakukan algoritme di
titik yang belum pernah ia lihat, rata-ratanya ia benar separuh kali, padahal label itu sebenarnya tertentu, sehingga $R^*=0$.

Pesannya bukan bahwa belajar mustahil, melainkan bahwa belajar \emph{tanpa asumsi} mustahil. Asumsi tentang distribusi, misalnya
bahwa fungsi targetnya mulus, bahwa ia hanya bergantung pada sedikit arah, atau bahwa ia hampir linear, adalah apa yang membuat
generalisasi mungkin. Inductive bias di \cref{ch:peta} adalah nama lain dari asumsi itu.

Kuliah merangkum ``perlombaan ilmiah untuk adaptivitas'' dengan beberapa laju konvergensi. Kalau fungsi target hanya diketahui
Lipschitz di $\R^d$, galat turun seperti $n^{-1/d}$. Untuk $d=10$, mengurangi galat setengahnya menuntut $2^{10}$, sekitar seribu
kali lebih banyak data. Inilah kutukan dimensi dalam bentuk yang paling telanjang. Kalau target punya $m$ turunan yang terbatas,
laju membaik menjadi $n^{-m/d}$, dan jaringan saraf bisa mencapai laju $n^{-m/k}$ kalau target sebenarnya hanya bergantung pada
proyeksi berdimensi $k$ dari input, sesuatu yang menurut naskah kuliah tidak bisa dicapai metode kernel. Algoritme yang baik adalah
algoritme yang menyesuaikan diri dengan struktur yang kebetulan ada, tanpa harus diberi tahu sebelumnya.

\section{Least squares di bawah kaca pembesar}

Kuliah memakai least squares sebagai laboratorium, karena semua efek penting, tawar-menawar bias dan varians, kutukan dimensi, dan
efek regularisasi, bisa dihitung secara eksak di sana. Modelnya linear dalam parameter, $f_{\bm\theta}(\vx)=\bm\theta^\top\phi(\vx)$, dengan
peta fitur $\phi$ yang boleh nonlinear dalam $\vx$. Estimator OLS $\hat{\bm\theta}=(\bm\Phi^\top\bm\Phi)^{-1}\bm\Phi^\top\vy$ diturunkan di
\cref{ch:data}, dan prediksinya adalah proyeksi ortogonal $\vy$ ke ruang kolom matriks fitur $\bm\Phi$.

Analisisnya paling bersih dalam \emph{fixed design}: input dianggap tetap, misalnya titik-titik grid, dan hanya output yang acak,
$y_i=\bm\theta^{*\top}\phi(\vx_i)+\epsilon_i$ dengan noise independen berrata-rata nol dan varians $\sigma^2$. Risiko diukur pada
input yang sama. Dengan $\hat{\bm\Sigma}=\tfrac1n\bm\Phi^\top\bm\Phi$, galat berlebih sebuah parameter adalah
$(\bm\theta-\bm\theta^*)^\top\hat{\bm\Sigma}(\bm\theta-\bm\theta^*)$. OLS tak bias, $\E[\hat{\bm\theta}]=\bm\theta^*$, dan kovariansnya
$\tfrac{\sigma^2}n\hat{\bm\Sigma}^{-1}$. Galat berlebih harapannya adalah jejak kovarians itu ditimbang dengan $\hat{\bm\Sigma}$:
\begin{equation}
\E\big[R(\hat{\bm\theta})\big]-R^*=\mathrm{tr}\Big(\hat{\bm\Sigma}\cdot\frac{\sigma^2}n\hat{\bm\Sigma}^{-1}\Big)=\frac{\sigma^2d}n .
\end{equation}
Rumus sesederhana ini memuat banyak pelajaran. Galat turun seperti $1/n$, laju terbaik yang mungkin. Tetapi ia sebanding dengan
jumlah parameter $d$. Dengan $d=50$, $n=1000$, dan $\sigma^2=1$, galat berlebihnya $0{,}05$. Dengan $d=500$, sepuluh kali lebih besar.
Kalau $d$ mendekati $n$, model bisa menghafal data latih, dan untuk $d=n$ dengan matriks yang bisa dibalik, prediksinya tepat sama
dengan data, noise sekalipun.

\section{Ridge dan derajat kebebasan}

Ridge regression menambahkan $\lambda\|\bm\theta\|^2$, dengan solusi $\hat{\bm\theta}_\lambda=\tfrac1n(\hat{\bm\Sigma}+\lambda\bm I)^{-1}\bm\Phi^\top\vy$.
Naskah kuliah memberi galat berlebih harapannya secara eksak sebagai jumlah dua suku:
\begin{equation}
\E\big[R(\hat{\bm\theta}_\lambda)\big]-R^*=\lambda^2\,\bm\theta^{*\top}(\hat{\bm\Sigma}+\lambda\bm I)^{-2}\hat{\bm\Sigma}\,\bm\theta^*
+\frac{\sigma^2}n\,\mathrm{tr}\big(\hat{\bm\Sigma}^2(\hat{\bm\Sigma}+\lambda\bm I)^{-2}\big).
\end{equation}
Suku pertama adalah bias kuadrat: naik dengan $\lambda$, karena parameter ditarik semakin kuat ke nol, dan tidak bergantung pada $n$.
Suku kedua adalah varians: turun dengan $\lambda$, mendekati $\sigma^2d/n$ ketika $\lambda$ menuju nol, dan menghilang ketika $\lambda$
besar. Jejak di suku kedua, $\mathrm{df}(\lambda)=\sum_i\lambda_i^2/(\lambda_i+\lambda)^2$ dengan $\lambda_i$ nilai eigen $\hat{\bm\Sigma}$,
disebut \istilah{derajat kebebasan} efektif.

Derajat kebebasan menjelaskan apa yang sebenarnya dilakukan regularisasi. Arah dengan nilai eigen jauh lebih besar dari $\lambda$
dihitung hampir penuh, arah dengan nilai eigen jauh lebih kecil hampir tidak dihitung. Untuk tiga arah dengan nilai eigen $4$, $1$,
dan $0{,}01$, dan $\lambda=0{,}1$, sumbangannya $16/16{,}81\approx0{,}95$, $1/1{,}21\approx0{,}83$, dan $0{,}0001/0{,}0121\approx0{,}01$,
sehingga $\mathrm{df}\approx1{,}79$. Model punya tiga parameter, tetapi variansnya seperti model dengan kurang dari dua parameter, karena
arah yang hampir tidak didukung data praktis dimatikan.

Dengan pilihan $\lambda$ yang tepat, naskah kuliah menunjukkan bahwa galat berlebih ridge bisa dibatasi oleh sebuah suku yang turun
seperti $1/\sqrt n$ dan sama sekali tidak bergantung pada $d$, hanya pada norma $\bm\theta^*$, besar fitur, dan $\sigma$. Bandingkan dengan
OLS: OLS turun lebih cepat, seperti $1/n$, tetapi dengan konstanta yang tumbuh dengan $d$. Ridge turun lebih lambat dengan konstanta
yang tidak bergantung pada $d$. Untuk model dengan banyak parameter dibanding data, keadaan normal di deep learning, yang kedua
jauh lebih menarik. Dalam praktik, $\lambda$ tidak dipilih dengan rumus karena $\sigma$ dan $\bm\theta^*$ tidak diketahui, melainkan
dengan validasi silang, yang menurut kuliah menikmati laju optimal yang sama.

\section{Kernel}

Bagaimana kalau peta fitur $\phi$ berdimensi besar sekali, bahkan tak hingga? Teorema representer menjawab bahwa untuk minimisasi
risiko empiris dengan penalti norma di ruang Hilbert, solusi optimalnya selalu bisa ditulis sebagai kombinasi fitur dari titik-titik
data, $\bm\theta=\sum_i\alpha_i\phi(\vx_i)$, apa pun loss-nya, asalkan penaltinya naik dengan norma. Pembuktiannya pendek. Pecah
$\bm\theta$ menjadi komponen di dalam rentang $\phi(\vx_i)$ dan komponen tegak lurus terhadapnya. Komponen tegak lurus tidak mengubah
satu pun prediksi di data latih, karena hasil kali dalamnya dengan setiap $\phi(\vx_i)$ nol, tetapi menambah norma. Jadi solusi
optimal tidak punya komponen tegak lurus.

Akibatnya, semua yang dibutuhkan hanya hasil kali dalam antar fitur, $k(\vx,\vx')=\phi(\vx)^\top\phi(\vx')$, yang disusun dalam matriks
kernel $\bm K$. Masalah di ruang berdimensi tak hingga menjadi masalah dengan $n$ variabel. Untuk loss kuadrat, kernel ridge regression
memberi $\bm\alpha=(\bm K+n\lambda\bm I)^{-1}\vy$ dan prediksi $\hat f(\vx)=\sum_i\alpha_ik(\vx_i,\vx)$, bentuk yang sama dengan rata-rata
posterior Gaussian process di \cref{ch:mlat}. Teorema Aronszajn menjamin bahwa setiap fungsi $k$ yang selalu menghasilkan matriks
semidefinit positif adalah hasil kali dalam di suatu ruang fitur, sehingga kita cukup merancang kernel, tanpa pernah menulis $\phi$.
Inilah \istilah{kernel trick}.

Contoh kecil memperlihatkan fitur tersembunyi di balik kernel. Untuk $\vx=(x_1,x_2)$, kernel polinomial $k(\vx,\vx')=(\vx^\top\vx')^2$
bisa diuraikan menjadi $(x_1x_1'+x_2x_2')^2=x_1^2x_1'^2+2x_1x_2x_1'x_2'+x_2^2x_2'^2$, yaitu hasil kali dalam dari fitur
$\phi(\vx)=(x_1^2,\ \sqrt2x_1x_2,\ x_2^2)$. Menghitung kernel hanya butuh satu hasil kali dalam dua dimensi lalu dikuadratkan,
sedangkan fiturnya tiga dimensi. Untuk pangkat dan dimensi yang lebih tinggi, selisihnya meledak. Kernel Gauss,
$k(\vx,\vx')=\exp(-\|\vx-\vx'\|^2)$, berkorespondensi dengan ruang fitur berdimensi tak hingga, dan tetap dihitung dengan satu jarak.

\section{Di luar kelas}

Teori di bab ini terdengar jauh dari praktik, tetapi ia menjawab pertanyaan yang muncul di seluruh buku. Mengapa ukuran data penting
untuk model pengganti simulasi berdimensi tinggi? Laju $n^{-1/d}$ menjawabnya, dan juga menjelaskan mengapa struktur fisika yang
mengurangi dimensi efektif begitu berharga (\cref{ch:pinn,ch:operator}). Mengapa jaringan saraf dengan parameter jauh lebih banyak dari
data tetap bisa generalisasi? Derajat kebebasan efektif dan regularisasi implisit memberi sebagian jawabannya, dan double descent di
\cref{ch:data} adalah kepingan lainnya. Dan teorema no free lunch adalah pengingat yang rendah hati: setiap keberhasilan sebuah metode
adalah keberhasilan asumsinya.

\sumberbab{Bab ini mengikuti naskah kuliah Statistical Learning Theory and Applications (pembelajaran terawasi dan risiko, minimisasi
risiko empiris, konsistensi dan risiko minimax, teorema no free lunch, analisis least squares dan ridge dalam fixed design, teorema
representer dan kernel). Rujukan utamanya adalah \citet{bach2024}.}
